\subsection{Aprendizaje por refuerzo para mejorar las respuestas}
Frente a esa limitación, el aprendizaje por refuerzo ofrece una forma de orientar el entrenamiento según la calidad de las respuestas generadas, en lugar de depender únicamente de la reproducción de ejemplos. Para ello, el modelo puede producir varias respuestas ante una misma consulta y recibir una recompensa que indique cuáles se acercan mejor a los objetivos definidos para el dominio \cite{ouyang2022instructgpt}. Es aquí donde entra en juego GRPO, que utiliza las recompensas de un grupo de respuestas a una misma pregunta para guiar la optimización y evita entrenar un modelo de valor independiente \cite{shao2024deepseekmath}. Su papel cobró relevancia en tareas con respuestas finales verificables, como los problemas matemáticos con resultados determinados, donde es posible comprobar automáticamente si la solución es correcta. DeepSeekMath introdujo este método para mejorar el razonamiento matemático y, en la evaluación sin herramientas externas, el entrenamiento con GRPO elevó la exactitud del modelo de 46,8\,\% a 51,7\,\% en MATH y de 82,9\,\% a 88,2\,\% en GSM8K, respecto del modelo previamente ajustado mediante entrenamiento supervisado \cite{shao2024deepseekmath}. Por su parte, el trabajo de DeepSeek-R1 amplió esta línea de investigación al estudiar cómo el entrenamiento con GRPO puede desarrollar capacidades de razonamiento a partir de recompensas basadas en reglas. En particular, DeepSeek-R1-Zero se entrenó directamente desde un modelo base, sin una etapa inicial de ajuste fino supervisado, utilizando recompensas de exactitud para verificar respuestas matemáticas con resultados determinados y soluciones de programación mediante casos de prueba, junto con recompensas de formato para distinguir el razonamiento de la respuesta final. Durante ese entrenamiento, los autores observaron comportamientos de revisión de las soluciones y exploración de estrategias alternativas, aunque también identificaron dificultades de legibilidad y mezcla de idiomas. Para abordar esas limitaciones, DeepSeek-R1 incorporó ejemplos supervisados iniciales y un proceso de entrenamiento en varias etapas que combina ajuste supervisado y aprendizaje por refuerzo \cite{deepseek2025r1}.
